三份材料摊在桌上：一篇物理论文说要``极小化变分自由能''，一篇机器学习论文说要``最大化 ELBO''，一份采样代码的注释说``温度越高越随机''。三句话用了三套词汇，你想知道它们是不是同一件事，还是各领域各自的黑话。答案是前两句严格是同一个式子，第三句里的温度就是那个式子中的一个系数。这一章要做的，就是把这套翻译关系建立起来，同时划清哪些对应是恒等式、哪些只是打比方。

这一章处理的对象是把概率写成 \(p(x)\propto e^{-E(x)}\) 的模型：Ising 模型、Boltzmann 机、受限玻尔兹曼机，以及任何一个只定义了未归一化打分的能量模型。它们的共同麻烦是配分函数——归一化常数是对指数多个状态求和，既算不出来也没法直接绕开。统计物理跟这个麻烦打了一百多年交道，攒下的三样工具正好可以整批搬过来：自由能把配分函数变成可求导的标量，平均场近似用独立假设把多体问题拆成单体自洽方程，相变则标出这些近似何时失效。

这一章要立住的判断有四条。第一，自由能把两种互相冲突的诉求写进一条式子：\(F=\E[E]-TS\)，能量项想让分布塌到最低处，熵项想让它摊平，温度是两者之间的兑换率——低温偏向能量，高温偏向熵。退火、采样温度、softmax 里的温度，调的都是这个兑换率。第二，物理里的变分自由能与机器学习里的证据下界是同一个量的两种写法，把能量取成负对数联合似然，配分函数就是证据，自由能的下界就是 ELBO，两边的间隙是同一个 \(\KL(q\|p)\)。第三，平均场近似的实质是假设变量之间独立，于是联合分布被拆成边缘分布的乘积；它系统性低估涨落，这与反向 KL 的模式寻找是同一件事的两种说法。第四，相变是集体现象：参数连续变化时宏观量突然换了性质，机制在于个体通过有效场互相强化。它与优化里``训练突然变好''或``突然崩''的相似之处是启发性的类比，不是定理。

\subsection{怎么读这一章}

三节各交一件可以带走的东西。第一节交一条恒等式 \(F[q]=T\KL(q\|p)+F\)，读懂它就同时读懂了 Boltzmann 分布为什么是平衡态、以及 ELBO 是从哪里来的。第二节交一个自洽方程 \(m=\tanh\bigl((Jm+h)/T\bigr)\) 和一条命题（乘积族里的方差不超过真方差），前者告诉你平均场怎么算，后者告诉你它一定往哪个方向错。第三节交一张分岔图和一份判据清单，用来判断你看到的陡变够不够格叫相变。

读的时候有两个开关要随时拨。一个是温度：物理文献带着 \(T\) 和 \(k_B\)，机器学习文献通常直接取 \(T=1\)，看到公式对不上时先检查是不是这个约定的差别。另一个是对象：同一个数学形式，在物理里描述真实粒子的平衡分布，在变分推断里描述一个计算上的近似，用错了指代就会把``粒子确实这样运动''和``我们只好这样近似''混为一谈。遇到深度学习的类比时，固定问三件事——对象是不是一个平衡分布，有没有明确的极限过程，序参量能不能写下来。

\subsection{本章篇目}

以下三篇按依赖顺序排列，也可以从当前最困惑的那一篇进入，再沿篇内链接回补。

\begin{itemize}
\tightlist
\item \hyperref[sec:13-Machine-Learning/13-Statistical-Physics-and-Free-Energy/01]{``自由能把配分函数变成热力学量''}；
\item \hyperref[sec:13-Machine-Learning/13-Statistical-Physics-and-Free-Energy/02]{``平均场近似用独立性换可计算性''}；
\item \hyperref[sec:13-Machine-Learning/13-Statistical-Physics-and-Free-Energy/03]{``相变与涨落：平均场在临界点附近失效''}。
\end{itemize}

配分函数本身的角色在\hyperref[sec:13-Machine-Learning/08-Sampling-and-Inference/06]{``配分函数把能量变成全局概率''}里已经交代过，变分近似的概率语言见\hyperref[sec:13-Machine-Learning/08-Sampling-and-Inference/03]{``变分推断把后验近似变成优化''}；这一章从那两处接手，换成物理的记法重讲一遍，好处是能量、温度、涨落这些量都有了明确的位置，坏处是词汇听起来比实际内容更玄。真正需要记住的分界线只有一条：配分函数、自由能、ELBO 之间是可以逐项代入验证的恒等式，而能量地形、相变、玻璃态这些说法在机器学习里目前多数只是有启发性的比喻，用它们提出猜想可以，拿它们当证据不行。
